\begin{abstract}
Persistent homology is increasingly computed on data that change over time, yet for general complexes of dimension two or more every known method for updating it after the insertion or deletion of a single simplex takes time polynomial in the number~$n$ of simplices in the worst case.
We show that this is no accident: for every $p\ge1$, maintaining the Betti number $\beta_p$ is equivalent, up to constant factors, to maintaining the rank of a sparse matrix.
Precisely, over the two-element field, the following two problems reduce to each other with constant overhead per operation: maintaining the first Betti number of a $2$-complex under simplex insertions and deletions, and maintaining the rank of a sparse matrix under entry updates; over every field, the same holds for matrices with entries in $\{0,\pm1\}$.
The key ingredient is a gadget, adapted from a static construction of Edelsbrunner and Parsa, in which switching a row or column off or on costs one simplex update and every entry update costs a constant number.
The switches alone imply that, under the Online Matrix--Vector multiplication (\OMv) conjecture, no data structure with polynomial preprocessing time can maintain a Betti number $\beta_p$, $p\ge1$, of a subcomplex of a fixed complex with $n$ simplices in amortized time $O(n^{1/2-\eps})$ per update and $O(n^{1-\eps})$ per query, for any $\eps>0$.
This holds over every field, and hence also for persistent Betti numbers under function-value updates and for online zigzag persistence; it survives the restriction to complexes of bounded degree, to $2$-complexes linearly embedded in~$\R^4$, and to Vietoris--Rips complexes of changing metric spaces and of changing point sets in~$\R^{O(\log n)}$, and it applies even to deciding whether the first homology group vanishes.
In contrast, all Betti numbers of subcomplexes of any fixed complex linearly embedded in~$\R^3$ can be maintained in $O(\log^2 n)$ amortized time per update after polynomial-time preprocessing.
Maintaining Betti numbers is thus easy for complexes in~$\R^3$, whereas for $2$-complexes in~$\R^4$ it is as hard as maintaining the rank of a sparse $0/1$ matrix, and over the two-element field equivalent to it.
\end{abstract}
